\section{State of the Modern World Address}

\textbf{Vietnam.} Let's say I hate Vietnam. Consider that more than 50,000 men of my generation died there, and more than that came home all messed up. People who have been everywhere, boast about their Vietnam vacation. Don't know why: all they can do there is visit a few Buddhist temples and the Nike factory.

The US Army was defeated there, but nothing bad happened. Vietnam has posed no threat to the USA. The so-called ``Domino Theory'', in which the surrounding countries were doomed to go communist, never actually happened. Yet the architects of that foreign policy never had to pay a price. In what other field can you be so wrong, yet keep your job?

I had to empathize with those scheduled to take their draft physical. Some got messed up on drugs: a week of meth and marijuana. Don't know if it was ever effective. Others told me they would pretend to be gay (the army did not want them then). I never figured out how you do that well enough to fool the doctors. I later found out that the physicians would look for a ``distended rectum'', whatever that is. So perhaps there was no pretense involved.

\textbf{Hindu Nationalism}. It may be tempting to conclude that the Hindu prayers to elect \textbf{Donald Trump}\footnote{\url{http://bigstory.ap.org/article/73ce604e1e2d40859fefab04f1e02d81/hindu-group-india-asks-gods-help-trump-win-us-election}} were more effective than the Christian prayers to elect Cruz, but it doesn't answer the why question.

\textbf{Narendra Modi}, the prime minister of India and a Trump-like figure, is commonly viewed as an anti-Muslim Hindu nationalist. The Hindu nationalists I've spoken to are convinced that Obama is a crypto-Muslim. They envision an international troika of Trump, Modi, and Putin opposed to both China and Pakistan. Let's call this the fifth political theory.

\textbf{The Wheel of Fortune}. We send our manufacturing to China. They take the money and build factories in Guatemala. The Guatemalans cannot live on the Chinese salary, so they come to the USA to work.

\textbf{The Shop around the Corner}. I recently watched this 1940 film on PBS\footnote{\url{http://www.imdb.com/title/tt0033045/}} with a young man, of the same age as most Gornahoor readers. I thought he would find it too tame and corny. Au contraire, he was utterly transfixed. Suddenly it dawned on him and he asked me, ``Is that what it was like before diversity?''

They even said, ``Merry Christmas'', he pointed out. The irony is that White Nationalists, who presumably would have preferred what the movie portrayed, would have had no role to play in that era. That movement is dominated by atheists, heathens, and inverts, the very beneficiaries of diversity.

\textbf{Party of Ideas}. Soi-disant American conservatives are always talking about being the ``party of ideas''. I'm afraid that liberalism is the party of ideas. Conservatism is the party of Tradition, Family, Property, Custom, Religion. It would be impossible to go through life having to intellectually justify every decision you make.

Another way to put it is that liberalism is nominalism, i.e., its ideas have no independent reality. Instead, they are artifacts, conventional, and factitious. Nominalism is opposed to realism, i.e., the eternal ideas that exist in God's mind and are then reflected in a healthy society.

\textbf{Corporate Power}. The socialist Bernie Sanders is always complaining about the influence of corporate power and money over politics. Has he ever complained about how corporate pressure has induced states like Indiana, North Carolina, and Georgia to alter their democratically enacted laws? This is not a matter of whether the laws are wise or not, but rather a question of who exactly legislates in this country. Why has no one ever asked Bernie about this?

\textbf{Overturning an election}. Yesterday, Brett Stephens of the WSJ announced on CNN that Donald Trump would have to be resoundingly defeated in order to teach the Republicans never to select such a candidate again. In today's follow up, Bill Kristol announced that there would be an independent candidate whose purpose would be to defeat Trump. (That candidate could never actually win.) Again, the issue is not whether Trump is a good candidate or not, just whether the voters have the right to elect him or not. Who exactly is this, obviously well-funded, group that determines which candidates are allowed to win an election?

\textbf{Building the Wall.} The Republican speaker of the house is married to a successful lobbyist for the Democrat party. She is worth millions. They live in a mansion in Wisconsin, surrounded by a wall built, presumably, to discourage interlopers from entering illegally. Common sense. So if a piece of legislation comes up that is important to the wife, what will Ryan do? Common sense.

\textbf{Fr. Charles Coughlin}. I recently found out from a relative that my grandfather was a big fan of Fr. \textbf{Charles Coughlin\footnote{\url{https://archive.org/details/Father\_Coughlin}} in the 1930s}. Coughlin was silenced by FDR and his journal could not be mailed. Of course, Coughlin was denounced as a Fascist, the easiest way to discount him, despite his criticism of the national socialists. His point was this: from a Christian perspective, the communists were more dangerous than the Fascists. Nevertheless, the USA sided with the communists. The talks are interesting for their historical interests. Some issue have never been resolved. Also, many of the laws under the Soviets are similar to laws in effect today in the USA.

\textbf{Imitation}. I was listening to an on-line podcast recently that sounded very much like a theme on Gornahoor. As it went on, it became too close to ideas written about here to be a coincidence. Actually, it sounded as if the man was simply reading the text.

\textbf{Paying off student loans}. It seems that more women are paying for their education by hooking onto sugar daddies\footnote{\url{http://www.zerohedge.com/news/2016-05-30/things-are-thriving-modern-hooker-economy}}. For some, this is morally dubious, although usually for the wrong reasons. Prostitution in itself is not a moral issue. What is immoral is the underlying fornication (or adultery if one of the parties is married). Whether or not money is exchanged does not alter the moral significance at all. Postscript: At one time, a man would have been praised for marrying a prostitute, since he took her from a life of fornication. Saving souls was considered praiseworthy.

\textbf{Milo Yiannopoulos}. I just found out about this fellow (yes, I am slow to keep up with popular culture). I watched a couple of his campus visits and debates. He is a young gay man who boasts of his talent for fellatio. Apparently he is the new face of ``conservative and catholic''. Somehow, he is my spiritual brother, although I don't believe he meets Evola's standard of what was considered healthy and normal.

\textbf{Diabolical Narcissism}. Although Milo is all over youtube, the uber-Catholic \textbf{Ann Barnhardt} has been banned from the site\footnote{\url{http://www.barnhardt.biz/2016/05/24/barnhardts-latest-video-banned-by-youtube-for-nudity/}}. Her latest video lecture is on the topic of ``Diabolical narcissism''. It starts off as an interesting idea combining a spiritual disease (diabolical) with its psychological manifestation (narcissism), and identifying its manifestation in different realms. Ultimately, it falls apart since it tries to account for too much; a cause that explains everything, explains nothing. For example, even Kabbalah is a manifestation. That would make Valentin Tomberg and Wolfgang Smith diabolical narcissists, which just does not seem likely.

It is worth a casual listen, nevertheless, if for no other reason than the occasional Barnhardt irreverence, such as Joe Biden being ``borderline retarded'' or the average IQ of the ``rap culture''. A debate between Milo and Ann would be something to look forward to.

\textbf{The Kali Yuga Is Everywhere}. In an essay published over a century ago, \textbf{D T Suzuki} laments the ``superficiality of our materialistic, industrial, culture'' in Japan, so he recommended Swedenborg to reinvigorate the Japanese religious spirit. In the west, many resent that same superficiality and turn to Buddhism as the remedy. Think on this and you may immediately grasp the irony: if you can't see the depth in your own tradition, you won't find it elsewhere.

\textbf{How to Attract Adherents}. Religious groups are always looking both to attract converts while preventing apostasy. For example, a Christian Science group was debating that issue. They wondered why healings were no longer happening today. Duh. If healings resulted from Christian Science treatments, people would be flocking to the Church.

The Orthodox have a different set of problems, judging from a discussion I stumbled upon. The parishes are too ethnic is one problem. Another issue is that all the converts from Protestantism are annoying, that is, anyone who actually believes in the dogmas and is not looking for a social club. I think the consensus conclusion was that they should be more gay friendly and concerned about world hunger.

\textbf{Sexually interesting}. I lost the link, but a British woman was boasting how all the other women at a party thought she was so interesting because of her extensive sexual experience. In saner times, to be interesting you would have needed the ability to read cuneiform tablets, or collected matchbooks from around the world, or perhaps had a unicorn ranch in some far off land.

\textbf{Refined Taste Buds}. In one of the stranger news stories of the past week, a woman claimed that there was semen in the sandwich she got from a fast food restaurant. I wondered how she knew. I felt vulnerable, since, unlike her, I would not have recognized the taste. That is why I never order anything with ranch dressing anymore (unless Milo comes with me).

\textbf{All's fair in War}. Once again the debate about whether the atomic bomb attacks against Japan has started: were they morally or militarily justified? The first question is easy: it is immoral to directly target civilians with no military purpose. The intent was to terrify the Japanese population. In a Machiavellian sense, the question is meaningless since the issue is the projection of power.

The answer, then, to the second question is it depends: if the goal was unconditional surrender, the bombing was successful. Since conditional surrender terms were previously rejected by the USA, it was the best option. The military defeat of Carthage was insufficient to the Romans; the entire city had to be destroyed. Machiavelli recommended that an enemy either had to be completely destroyed or else become one of the victor's allies. The Japanese chose the second option.


\flrightit{Posted on 2016-05-31 by Cologero}

\begin{center}* * *\end{center}

\begin{footnotesize}
\begin{sffamily}
\texttt{Wayne Ferguson on 2016-05-31 at 09:05 said:} 

That is why I never order anything with ranch dressing anymore (unless Milo comes with me).'' Not sure If that's a double entendre or a Freudian slip, but it's funny... :-)

\hfill

\texttt{Mark Citadel on 2016-05-31 at 10:04 said:}

I'm sorry you had to hear about Milo Yiannopoulos. A lot of people on the `alt-right' are fixated with him, but I find the man odious in the extreme. He's a contrarian in every meaning of the word. If Traditionalism dominated, he'd be a flaming liberal.

The Orthodox Church in the West has done a poor job of keeping Liberal influences out. I wonder if there is any iron cloister that could protect people from the spirit of the age, should it enter the vortex of insanity that circles Western Europe and the United States. I go to a Greek Orthodox Church (no Russian ones nearby), and my priest has not yet said anything I found objectionable, but always I will defer to priests of good character in Russia over anyone in the West who might be compromised by silly ideas about marriage and the like.

\hfill

\texttt{Cologero on 2016-06-04 at 14:05 said:}

The Freudian slip is in your own mind Wayne, but we try to appeal to all types.

\hfill

\texttt{nathan on 2022-12-24 at 01:04 said:}

I'm interested, if someone has a reference or can possibly provide one, regarding that particular essay published by D T Suzuki lamenting the ``superificiality of ... materialistic, industrial culture `in Japan.' ''

Would like to have this on record if at all possible.

\hfill

\texttt{Mike on 2023-05-31 at 11:34 said:}

@nathan

Swedenborg: Buddha of the North, Chapter 1

\end{sffamily}
\end{footnotesize}

